\chapter{Valuation and Financial Statements}\label{ch:iv2:valuation-and-financial-statements}

\begin{ivtrains}
The fundamental analyst's interview: following one transaction through three statements, building a quick valuation by hand, saying what a multiple assumes, reading earnings for what they hide, and defending a position when the numbers move. Each stem gives every number you need. Work on paper, say what each step means, and finish with what you would do.
\end{ivtrains}

\begin{question}[$\star$ \iqroles{Analyst} \iqfirm{hedge fund, bank} \iqlevel{junior}]\label{iq:iv2:valuation-and-financial-statements:1}
A company's depreciation rises by 10 for the year, and the tax rate is 25\%. Depreciation is deductible in the same amount for tax, and nothing else changes. Walk through the income statement, the cash flow statement and the balance sheet. By how much do net income, operating cash flow, cash, fixed assets and retained earnings change?
\end{question}

\begin{question}[$\star$ \iqroles{Analyst} \iqfirm{hedge fund, bank} \iqlevel{junior}]\label{iq:iv2:valuation-and-financial-statements:2}
A new company buys goods for 60 on credit and sells them for 100 on credit. At year end it has neither paid its supplier nor collected from its customer. It owes 25\% tax on its profit and pays it in cash during the year. What are its net income and its operating cash flow? Show that the balance sheet balances.
\end{question}

\begin{question}[$\star$ \iqroles{Analyst, Trader} \iqfirm{hedge fund, bank} \iqlevel{junior}]\label{iq:iv2:valuation-and-financial-statements:3}
A company has a market capitalisation of 1\,000, debt of 300, cash of 100, preferred stock of 50 and minority interest of 30, and EBITDA, its earnings before interest, tax, depreciation and amortisation, is 160. What are its enterprise value and its EV to EBITDA multiple? Does the enterprise value change if it borrows 100 and holds the cash, pays a dividend of 50, or spends 100 of cash on a share buyback?
\end{question}

\begin{question}[$\star$ \iqroles{Analyst} \iqfirm{hedge fund, bank} \iqlevel{junior}]\label{iq:iv2:valuation-and-financial-statements:4}
A company has operating profit of 200, a tax rate of 25\%, depreciation and amortisation of 40, capital expenditure of 70, and working capital that rises by 20. It pays interest of 20 and borrows 10 more on a net basis. What are its free cash flow to the firm and its free cash flow to equity? Which of them do you discount at the weighted average cost of capital?
\end{question}

\begin{question}[$\star$ \iqroles{Analyst, QR} \iqfirm{hedge fund, bank} \iqlevel{junior}]\label{iq:iv2:valuation-and-financial-statements:5}
A company will produce free cash flow of 100 next year, growing 3\% a year forever, and its cost of capital is 9\%. What is it worth? What is it worth if growth is 4\%, or if the cost of capital is 8\%? What changes if the 100 is this year's cash flow, not next year's?
\end{question}

\begin{question}[$\star$ \iqroles{Analyst} \iqfirm{hedge fund, bank} \iqlevel{junior}]\label{iq:iv2:valuation-and-financial-statements:6}
A stock trades at 25, with 100 shares outstanding. There are also 10 options with a strike of 10, 2 restricted shares that have not yet vested, and a convertible bond with a face value of 200 that converts at 20 per share. How many diluted shares are there, and what is the diluted equity value? What do you do with the convertible's 200 when you go from equity value to enterprise value?
\end{question}

\begin{question}[$\star$ \iqroles{Analyst, Trader} \iqfirm{hedge fund, bank} \iqlevel{junior}]\label{iq:iv2:valuation-and-financial-statements:7}
A company has revenue of 100, variable costs of 20\% of revenue, fixed costs of 60 and interest of 10. If revenue rises by 10\%, by what percentage do operating profit and pre-tax profit rise? At what revenue does operating profit reach zero, and at what revenue does pre-tax profit?
\end{question}

\begin{question}[$\star\star$ \iqroles{Analyst, QR} \iqfirm{hedge fund, bank} \iqlevel{junior}]\label{iq:iv2:valuation-and-financial-statements:8}
A company has equity worth 800 and debt worth 200. Its equity beta is 1.2, the risk-free rate is 4\%, the equity risk premium is 5\%, its cost of debt is 5\% and its tax rate is 25\%. What is its weighted average cost of capital? It then moves to equal amounts of debt and equity by value, and its cost of debt rises to 6\%. Removing and restoring the leverage with the formula that assumes debt has a beta of zero, what are its new equity beta and new weighted average cost of capital?
\end{question}

\begin{question}[$\star\star$ \iqroles{Analyst, QR} \iqfirm{hedge fund, bank} \iqlevel{junior, experienced}]\label{iq:iv2:valuation-and-financial-statements:9}
A company with no debt earns 100 after tax every year forever and pays it all out. Its cost of capital is 10\%, and the tax rate is 25\%. It borrows 500 forever at 5\%, risk-free, and uses the money to buy back shares. Assume that interest is tax-deductible and that there are no costs of financial distress. What is the company worth now, what is its equity worth, what is its cost of equity, and what is its weighted average cost of capital? Why is that not about 7\%?
\end{question}

\begin{question}[$\star\star$ \iqroles{Analyst} \iqfirm{hedge fund, bank} \iqlevel{junior, experienced}]\label{iq:iv2:valuation-and-financial-statements:10}
Free cash flow is 100 next year and grows 10\% a year for the following four years. After year 5 it grows 3\% a year forever. The cost of capital is 9\%, and each year's cash arrives at the end of the year. Value the company, and say what share of the value comes from the terminal value. What does the value become if terminal growth is 4\%? How does the value change if cash arrives in the middle of each year?
\end{question}

\begin{question}[$\star\star$ \iqroles{Analyst} \iqfirm{hedge fund, bank} \iqlevel{junior, experienced}]\label{iq:iv2:valuation-and-financial-statements:11}
In a valuation model, year 5 EBITDA is 100, year 5 free cash flow is 55 and the cost of capital is 9\%. A colleague uses a terminal value of 10 times EBITDA. What growth rate after year 5 does that imply? What multiple would 2\% growth imply? A colleague proposes 18 times: what growth does that require, and what do you say?
\end{question}

\begin{question}[$\star\star$ \iqroles{Analyst, QR} \iqfirm{hedge fund, bank} \iqlevel{junior, experienced}]\label{iq:iv2:valuation-and-financial-statements:12}
A company earns 100 after tax on its operations next year, and it has a 10\% cost of capital. It will grow 5\% a year forever by reinvesting at a constant return on invested capital. What is it worth if that return is 5\%, 10\% or 20\%? What is it worth with no growth? Derive the formula you use. When does growth add value?
\end{question}

\begin{question}[$\star\star$ \iqroles{Analyst} \iqfirm{hedge fund, bank} \iqlevel{junior, experienced}]\label{iq:iv2:valuation-and-financial-statements:13}
Two companies have the same operations: operating profit of 100, depreciation of 20, a tax rate of 25\% and an enterprise value of 1\,200. Company $A$ has no debt. Company $B$ has debt of 600 at 5\%. What are the EV to EBITDA multiple and the price to earnings multiple for each? Why does the second differ? What would it be for $B$ if the debt cost 12\%?
\end{question}

\begin{question}[$\star\star$ \iqroles{Analyst} \iqfirm{hedge fund, bank} \iqlevel{junior, experienced}]\label{iq:iv2:valuation-and-financial-statements:14}
A company with no tax has revenue less other costs of 200 and spends 100 on developing a product. Under one policy it expenses the 100 at once. Under the other it capitalises it and amortises it evenly over 5 years, starting this year. For this first year, what are EBITDA, operating profit, operating cash flow, investing cash flow and free cash flow under each? What happens to the difference in operating profit if it spends 100 every year for many years, and if the spending grows 20\% a year?
\end{question}

\begin{question}[$\star\star$ \iqroles{Analyst} \iqfirm{hedge fund, bank} \iqlevel{junior}]\label{iq:iv2:valuation-and-financial-statements:15}
A company has revenue of 1\,000 and cost of goods sold of 600. It collects from customers after 45 days, holds inventory for 60 days and pays suppliers after 30 days. Take 365 days in a year, and base receivables on revenue, and inventory and payables on cost of goods sold. What are its cash conversion cycle and its net working capital? How much cash does 20\% growth absorb, if the days stay the same? How much does cutting inventory by 15 days release at the current size?
\end{question}

\begin{question}[$\star\star$ \iqroles{Analyst} \iqfirm{hedge fund, bank} \iqlevel{junior, experienced}]\label{iq:iv2:valuation-and-financial-statements:16}
A company reports free cash flow of 150, with share-based pay of 50 added back as a non-cash item. It has 100 shares at 30. What free cash flow yield does management quote, and what is the yield if you treat the share-based pay as a cost? How many shares would it have to buy back each year to keep the count flat? At 20 times free cash flow, how far apart are the two valuations?
\end{question}

\begin{question}[$\star\star$ \iqroles{Analyst} \iqfirm{hedge fund, bank} \iqlevel{junior, experienced}]\label{iq:iv2:valuation-and-financial-statements:17}
Company $A$ earns 100 after tax, with 100 shares at 20. It buys company $B$, which earns 20 after tax, for 300. There are no synergies, and the tax rate is 25\%. What is the new earnings per share if it pays in new shares? What is it if it borrows the 300 at 6\%? What rule decides each case? What happens with shares if the price is 450? If a deal adds to earnings per share, does that show that it creates value?
\end{question}

\begin{question}[$\star\star$ \iqroles{Analyst} \iqfirm{hedge fund, bank} \iqlevel{junior, experienced}]\label{iq:iv2:valuation-and-financial-statements:18}
A buyout fund pays 10 times EBITDA of 100, funded with debt of 600 and equity of 400. EBITDA grows 5\% a year for 5 years, and free cash flow repays 300 of debt over that time. It sells at 10 times EBITDA after 5 years. What are the equity value at exit, the multiple of money and the annual return? Where does the gain come from? What is the return if it sells at 9 times?
\end{question}

\begin{question}[$\star\star\star$ \iqroles{Analyst} \iqfirm{hedge fund, bank} \iqlevel{junior, experienced}]\label{iq:iv2:valuation-and-financial-statements:19}
A company with an enterprise value of 1\,000 and EBITDA of 150 rents a building for 100 a year, paid at year end, on a lease with 5 years left. It has been treating the rent as an operating cost. It now reports under a rule that puts the lease on the balance sheet at the present value of the payments, discounted at 6\%. What are the lease liability and the asset? What are the first year's depreciation, interest, EBITDA change and pre-tax profit change? What is the EV to EBITDA multiple after the change? What happens to the cash, and to the total expense over the 5 years?
\end{question}

\begin{question}[$\star\star\star$ \iqroles{Analyst} \iqfirm{hedge fund, bank} \iqlevel{junior, experienced}]\label{iq:iv2:valuation-and-financial-statements:20}
A company buys a machine for 100 and earns 60 a year before depreciation for 5 years. It depreciates the machine evenly over 5 years in its accounts, and over 2 years, evenly, for tax. The tax rate is 25\%. For each year, give net income, cash tax, the deferred tax balance and operating cash flow. What is the deferral worth at a discount rate of 8\%?
\end{question}

\begin{question}[$\star\star\star$ \iqroles{Analyst} \iqfirm{hedge fund, bank} \iqlevel{junior, experienced}]\label{iq:iv2:valuation-and-financial-statements:21}
A company pays 500 in cash for all of a target whose book equity is 200. The target's fixed assets are worth 100 more than their book value, and neither the write-up nor the goodwill is deductible for tax. The tax rate is 25\%. How much goodwill does the buyer record? If the write-up is depreciated over 10 years, how does that affect the buyer's net income each year? Later the goodwill is impaired by 100: what happens to net income, tax and cash?
\end{question}

\begin{question}[$\star\star\star$ \iqroles{Analyst, QR} \iqfirm{hedge fund, bank} \iqlevel{junior, experienced}]\label{iq:iv2:valuation-and-financial-statements:22}
A company has an enterprise value of 3\,000 and just produced free cash flow of 60. Its cost of capital is 8\%. What constant growth rate forever does the price imply? If instead growth is 3\% forever after 10 years, what growth rate over those 10 years does the price imply? What is the company worth if it grows at 3\% from today? How would you use these numbers?
\end{question}

\begin{question}[$\star\star\star$ \iqroles{Analyst} \iqfirm{hedge fund, bank} \iqlevel{junior, experienced}]\label{iq:iv2:valuation-and-financial-statements:23}
A retailer has 100 stores, each with revenue of 2 and a store-level margin of 15\%, and central costs of 12. Next year sales at existing stores rise 3\%, and 30\% of any change in their revenue falls through to store profit, because much of their cost is fixed. Central costs stay the same. It also opens 10 new stores in the middle of the year. A new store sells at an annual rate of 70\% of a mature store's current revenue, with a 5\% margin. Build next year's revenue and operating profit by hand. What is the growth in operating profit? What does one point of same-store sales do to profit, and what if same-store sales are flat?
\end{question}

\begin{question}[$\star\star\star\star$ \iqroles{Analyst} \iqfirm{hedge fund} \iqlevel{experienced}]\label{iq:iv2:valuation-and-financial-statements:24}
A company reports these figures for two years, and management calls the year ``a record''. The tax rate is 25\%. Assets were 1\,000 at the start of year 2 and 1\,250 at the end.
\begin{center}\small
\begin{tabular}{@{}lrr@{}}
 & Year 1 & Year 2 \\ \hline
Revenue & 1\,000 & 1\,200 \\
Cost of goods sold & 600 & 720 \\
Net income & 100 & 130 \\
Operating cash flow & 110 & 60 \\
Capital expenditure & 40 & 40 \\
Development costs capitalised & 20 & 70 \\
Receivables & 120 & 210 \\
Inventory & 100 & 170 \\
Deferred revenue & 80 & 60
\end{tabular}
\end{center}
Which lines worry you? Compute what you need to back that up. What did earnings grow by if development costs had stayed at Year 1's share of revenue? What innocent explanations would you test, and what would you ask management?
\end{question}

\begin{question}[$\star\star\star\star$ \iqroles{Analyst, Trader} \iqfirm{hedge fund} \iqlevel{experienced}]\label{iq:iv2:valuation-and-financial-statements:25}
You own a stock at 90, sized so that a fall to your bear case loses 1\% of your capital. You valued it at 150 if margins recover, and 60 if they do not, and you think recovery has a 50\% chance. The company then cuts guidance, and the stock falls to 63. You now value the two cases at 120 and 50. What recovery probability did the old price imply, and what does the new price imply? What probability do you need for a 15\% expected return from here? If you now believe 35\%, what is your expected return? How large would a fresh position be under the same 1\% rule, and what did you lose so far? Do you add, hold or cut, and what would change your mind?
\end{question}
